\subsection{QC-012: a complete three-copy degree-six nonic restriction}

\paragraph{Statement QC-012.}
At $(m,n)=(8,9)$, for
\[
                  d=6,\qquad \lambda=(14,14,14,12),
\]
the complete multiplicities are
\[
                      a_\lambda=q_\lambda=K_9(\lambda)=3.
                      \tag{QC-012.1}
\]
In particular this label gives no actual padded-permanent multiplicity
obstruction.  The determinant assertion follows from a three by three
restriction minor on explicit, genuine $9\times9$ pencils.  It does not
use a stabilizer occurrence bound or a conjectural canonical Foulkes map.

\paragraph{Three integral highest-weight polynomials.}
Introduce six independent symbolic vectors $u_0,\ldots,u_5$ with four
coordinates.  Write $[ijkl]$ for the determinant of these four vectors,
in the indicated order, and $[ijk]$ for their determinant in coordinates
$0,1,2$.  Define
\begin{align*}
 \Omega_0&=[0123]^4[0145]^4[2345]^4[024][135],\\
 \Omega_1&=[2345]^4[0145]^3[0123]^3[0135]^2[024][124],\\
 \Omega_2&=[2345]^3[0145]^3[0123]^2[1345]^2[0235][0234][012]^2.
\end{align*}
Every symbolic vector has degree nine in each expression.  There are
twelve four-coordinate brackets and two three-coordinate brackets,
counting repetitions.  The coordinate weight is therefore
$(14,14,14,12)$.  Each upper simple-root operator
$\sum_s (u_s)_i\partial_{(u_s)_{i+1}}$ annihilates every bracket: it
either differentiates a missing row or replaces a determinant row by a
row already present.  Consequently these are highest-weight symbols.

For an ordinary quaternary nonic $F$, set
\[
 P_i(F)=\left.
  \Omega_i(\partial_{u_0},\ldots,\partial_{u_5})
          \prod_{s=0}^5F(u_s)\right|_{u_0=\cdots=u_5=0}.
          \tag{QC-012.2}
\]
These are integral, degree-six coefficient polynomials.  The identical
six inputs automatically symmetrize the letters.  No nonvanishing or
independence of that symmetrization is assumed; the restriction matrix
below proves both for the three selected symbols.

\paragraph{Three actual determinant points.}
Let
\[
 M_1=\begin{pmatrix}
 x_3&x_0&x_1\\x_0&x_3&x_2\\-x_1&x_2&x_3
 \end{pmatrix},\qquad
 M_2=\begin{pmatrix}
 x_3&x_0&0\\0&x_3&x_1\\x_2&0&x_3
 \end{pmatrix}.
\]
Their determinants, checked directly by the six-term Leibniz formula,
are
\[
 H_1=x_3^3-x_0^2x_3+x_1^2x_3-x_2^2x_3,
 \qquad H_2=x_3^3+x_0x_1x_2.
\]
Put
\[
 F_1=(x_0x_1x_2)^2H_1,\qquad
 F_2=(x_0x_1x_2)^2H_2,\qquad
 F_3=F_1(x_3,x_1,x_2,x_0).
\]
For $i=1,2$, an explicit pencil for $F_i$ is
\[
       \operatorname{diag}(x_0,x_0,x_1,x_1,x_2,x_2,M_i).
\]
Substituting $x_0\leftrightarrow x_3$ in the first pencil gives the
third.  Thus all three points belong to the determinant image itself.

\paragraph{The exact restriction matrix.}
With the normalization in (QC-012.2), the reconstructed matrix modulo
$101$ is
\[
  \bigl(P_i(F_j)\bigr)_{0\leq i\leq2,\,1\leq j\leq3}
       =\begin{pmatrix}
          81&5&15\\22&95&1\\27&2&72
        \end{pmatrix},\qquad
       \det=46\pmod {101}.
       \tag{QC-012.3}
\]
The determinant is a nonzero integer modulo $101$, so it is nonzero
over $\mathbb C$.  Restriction of the three-dimensional highest-weight
space is therefore injective on every copy simultaneously.

\paragraph{Complete coefficient computation.}
For an exponent vector $\alpha$, write $u^\alpha$ for an ordinary
monomial, without divided powers.  The five exponents needed are
\[
 \begin{array}{c|ccccc}
       &A&B&C&D&E\\\hline
 \alpha&(2,2,2,3)&(4,2,2,1)&(2,4,2,1)&(2,2,4,1)&(3,3,3,0).
 \end{array}
\]
Here $F_1=A-B+C-D$ and $F_2=A+E$.  Coordinate weights force the
entire sixfold product contributing to $P_i(F_1)$ to have multiset
$A^3BCD$, and that contributing to $P_i(F_2)$ to have multiset
$A^4E^2$.  The checker enumerates all sixfold monomial multisets to
verify that no other one contributes.  Moving the coordinate
transposition in $F_3$ onto the symbol replaces each three-bracket by
the bracket in coordinates $1,2,3$, with total sign $+1$.  The required
multiset in $F_1^6$ is then $A^4CD$, whose coefficient has sign $-1$.
There are respectively $120$, $15$, and $30$ ordered assignments.

The sums of the corresponding symbolic coefficients, before their
ordinary differential factors, are
\[
 \begin{pmatrix}
 39&56&8\\48&54&14\\13&83&99
 \end{pmatrix}\pmod {101}.
\]
For a fixed ordered assignment $\alpha^{(0)},\ldots,\alpha^{(5)}$,
differentiation contributes exactly
$\prod_{s,k}\alpha^{(s)}_k!$.  These factors are respectively
$72$, $56$, and $36$ modulo $101$, independent of the ordering within
each multiset.  Multiplying the three columns by $72$, $56$, and
$-36$ gives exactly (QC-012.3).  No factorial normalization is left
implicit.

Every four-bracket is expanded by its $24$ signed permutations, and
every three-bracket by its six signed permutations.  Direct convolution
gives $282$, $1848$, and $10147$ nonzero modular monomials for the
second, third, and fourth powers of a generic four by four determinant.
In each symbol select the factors supported on $0145$, $0123$, and
$2345$.  The first two share letters $0,1$; after their exponents there
are matched, the required coefficient of the third is a dictionary
lookup.  All remaining factors are convolved explicitly.  For
$\Omega_2$ there are $940882$ such small-factor monomials, and
preconvolving the first two large factors gives $362712$ distinct
monomials in $2296$ groups indexed by the exponents of letters $0,1$.
This is an exact reassociation of a finite polynomial product.

The implementation encodes each exponent in five bits.  No exponent
exceeds nine, so addition has no digit carry.  For subtraction, it adds
$16$ separately in each digit before subtracting the used exponents;
the $16$ bit detects each negative residual.  The degrees of the
factors ensure that every intermediate digit remains between zero and
31, so neither subtraction nor the test mixes adjacent coordinates.
The two remaining coefficient-joining implementations use the same
explicit coefficient rule; one sums over a large-factor monomial first,
and the other first convolves the two large factors.

To reduce repeated work, the checker enumerates all $720$ letter
permutations and retains only exact positive automorphisms of each
displayed symbol.  It also uses coordinate permutations among the
three chosen bracket coordinates.  Their determinant signs cancel,
because there are twelve four-brackets and two three-brackets.
Each ordered target is assigned to a canonical representative, and its
multiplicity is retained.  Thus the quotienting is an explicitly
checked identity of coefficient sums, not a random sampling step.

\paragraph{Ambient count and conclusion.}
The same fresh checker computes the ambient multiplicity by coefficient
counting in
\[
 \prod_{\alpha_0+\alpha_1+\alpha_2+\alpha_3=9}
       (1-tx^\alpha)^{-1}.
\]
Its 24-term Weyl alternation at $(14,14,14,12)$ is exactly three.
The complete list of integer weight counts is included in the JSON
result.  Together with (QC-012.3) this gives $K_9(\lambda)=a_\lambda=3$.
QC-011 gives $q_\lambda=a_\lambda$ since $d=6\leq9$.
The four-variable reduction identifies these complete highest-weight
spaces with the four-row multiplicities in the main question.

\paragraph{Reproduction and scope.}
Run
\begin{quote}\ttfamily\small
python work/quaternary/scripts/check\_three\_nonic\_restrictions.py
\end{quote}
The checker uses only the Python standard library and reconstructs all
polynomials, bracket powers, ordered coefficient sums, and the ambient
count afresh.  It reads no stored matrix or pickle cache.  Its output
is \texttt{work/quaternary/results/three\_nonic\_restrictions.json}.
The exploratory scripts and their earlier dependent trial symbol are
retained separately; they are not prerequisites of this certificate.
This is an explicit small multiplicity-space restriction certificate.
It makes no assertion that the entire degree-six determinant ideal
vanishes, or that other degree-six labels have full determinant rank.
